Put \(m=\beta-1\), and index the coordinates of \(\mathsf H_\beta\) by \(0,\ldots,m\).  For every \(a=(a_0,\ldots,a_m)^{\mathsf T}\in\mathbb R^{m+1}\), the definition of \(\mathsf H_\beta\) gives
\[
a^{\mathsf T}\mathsf H_\beta a
=\sum_{j=0}^m\sum_{\ell=0}^m\frac{a_ja_\ell}{j+\ell+1}
=\int_0^1\left(\sum_{j=0}^m a_jx^j\right)^2\,dx.
\]
If \(a\ne0\), the polynomial in the last integrand is nonzero.  Its square is continuous, nonnegative, and strictly positive on a nonempty open subinterval of \([0,1]\), so the integral is strictly positive.  Thus \(\mathsf H_\beta\) is positive definite and hence invertible.

We next compute the squared distance of \(x^m\) from the polynomials of degree less than \(m\).  Define
\[
L_m(x)=\frac{1}{m!}\frac{d^m}{dx^m}\left\{x^m(x-1)^m\right\}.
\]
The coefficient of \(x^m\) in \(L_m\) is
\[
\frac{1}{m!}\frac{(2m)!}{m!}=\binom{2m}{m}.
\]
If \(h\) is any polynomial with \(\deg(h)<m\), integration by parts \(m\) times yields
\[
\int_0^1h(x)L_m(x)\,dx
=\frac{(-1)^m}{m!}\int_0^1h^{(m)}(x)x^m(x-1)^m\,dx
=0.
\]
Indeed, every boundary term vanishes because \(x^m(x-1)^m\) and each of its derivatives of order less than \(m\) vanish at both endpoints, while \(h^{(m)}=0\).  The same integration-by-parts identity with \(h=L_m\), now using \(L_m^{(m)}=m!\binom{2m}{m}\), gives
\[
\begin{aligned}
\int_0^1L_m(x)^2\,dx
&=\binom{2m}{m}\int_0^1x^m(1-x)^m\,dx\\
&=\binom{2m}{m}\frac{(m!)^2}{(2m+1)!}
=\frac{1}{2m+1}.
\end{aligned}
\]
For completeness, the middle integral follows without an external formula: if \(I_{a,b}=\int_0^1x^a(1-x)^b\,dx\), one integration by parts gives \(I_{a,b}=bI_{a+1,b-1}/(a+1)\); iterating from \((a,b)=(m,m)\) and using \(I_{2m,0}=1/(2m+1)\) yields \(I_{m,m}=(m!)^2/(2m+1)!\).

Consequently
\[
P_m(x)=\binom{2m}{m}^{-1}L_m(x)
\]
is monic of degree \(m\), is orthogonal in \(L^2[0,1]\) to every polynomial of degree less than \(m\), and satisfies
\[
\int_0^1P_m(x)^2\,dx
=\frac{1}{(2m+1)\binom{2m}{m}^{2}}.
\]
Since \(x^m-P_m(x)\) has degree less than \(m\), orthogonality shows, for every polynomial \(g\) of degree less than \(m\),
\[
\int_0^1\{x^m-g(x)\}^2\,dx
=\int_0^1P_m(x)^2\,dx
+\int_0^1\{x^m-P_m(x)-g(x)\}^2\,dx.
\]
It follows that the first term on the right is the minimum squared norm among all monic degree-\(m\) polynomials.

When \(m\ge1\), partition the positive-definite Gram matrix as
\[
\mathsf H_\beta=
\begin{pmatrix}
A&b\\
b^{\mathsf T}&c
\end{pmatrix},
\]
where \(A\) corresponds to \(1,x,\ldots,x^{m-1}\).  The preceding minimum and completion of the square give
\[
s:=c-b^{\mathsf T}A^{-1}b
=\min_{u\in\mathbb R^m}
\begin{pmatrix}u\\1\end{pmatrix}^{\!\mathsf T}
\mathsf H_\beta
\begin{pmatrix}u\\1\end{pmatrix}
=\frac{1}{(2m+1)\binom{2m}{m}^{2}}.
\]
Here \(A\) is invertible because it is a principal submatrix of the positive-definite matrix \(\mathsf H_\beta\), and \(s>0\) by the displayed minimum.  Direct multiplication, rather than an undeclared inverse formula, shows
\[
\mathsf H_\beta
\begin{pmatrix}-A^{-1}b\\1\end{pmatrix}
=
\begin{pmatrix}0\\s\end{pmatrix}.
\]
Therefore the last coordinate of \(\mathsf H_\beta^{-1}e_{m+1}\) is \(s^{-1}\), so
\[
(\mathsf H_\beta^{-1})_{\beta,\beta}
=s^{-1}
=(2m+1)\binom{2m}{m}^{2}
=(2\beta-1)\binom{2\beta-2}{\beta-1}^{\!2}.
\]
If \(m=0\), then \(\mathsf H_1=(1)\), and the same formula equals \(1\).  This proves both assertions for every integer \(\beta\ge1\).